\chapter{Introdução}

Desde que o primeiro veículo a combustão interna foi produzido lá no final do século XIX por Karl Benz, a ideia geral de automóvel não sofreu grandes alterações,
mantendo uma essencia quase completamente mecânica por mais de um século, usando o motor a combustão para propulção e as rodas e eixos para transmissão de potência
para o solo. Mas como destaca Hodel \cite{ref:hodel}, o que mais mudou nessas ultimas décadas foi o desenvolvimento de sistemas eletrônicos embarcados nos veículos,
saindo de um produto puramente mecânico para um produto que cada vez mais intregra sistemas computacionais e softwares.

Essa revolução tecnológica vem após a introdução das Unidades de Controle Eletrônico (ECUs) nos veículos. Conforme destaca Jubelli \cite{ref:jubelli}, tudo se
iniciou por volta de 1970 e evoluiu para os atuais sistemas de controle eletrônicos embarcados, onde cada carro pode possuir mais de 80 ECUs, responsáveis por
controlar e monitorar uma variedade de funções do automóvel, gerenciando sistema criticos como o motor, transmissão e segurança. O aumento exponencial da
complexidade empregada nesses sistemas impulsionou o surgimento de Sistemas de Transporte Inteligente (ITS), que buscam integrar comunicação e controle para
melhorar a segurança, eficiência e sustentabilidade do transporte.

Como parte da evolução dos ITS, hoje já falamos em Sistemas de Transporte Inteligente Cooperativo (C-ITS), onde a conciência cooperativa é concebida através da
tecnologia V2X (\textit{Vehicle-to-Everything}). Automóveis convencionais geralmente possuem sensores como radares e câmeras para monitorar o ambiente ao redor,
mas acabam sendo limitados pela linha de visão e alcance desses sensores. O V2X, por outro lado, permite trazer uma visão além da linha de visão do próprio veículo,
possibilitando a comunicação entre outros veículos, infraestrutura rodoviária e até mesmo pedestres, criando uma maior conciência sobre o ambiente próximo. Essa
comunicação cooperativa foi viabilizada pelas tecnologias DSRC ou C-V2X, que funcionam respectivamente pelos padrões da IEEE 802.11p e por redes celulares LTE e 5G,
permitindo a troca de informações em tempo real entre os participantes do sistema de transporte.

\abreviatura{C-ITS}{Sistema de Transporte Inteligente Cooperativo}
\abreviatura{DSRC}{\textit{Dedicated Short-Range Communications}}
\abreviatura{C-V2X}{\textit{Cellular Vehicle-to-Everything}}
\abreviatura{SUMO}{\textit{Simulation of Urban Mobility}}


\section{Justificativa}

Texto texto texto texto texto texto texto texto texto texto texto texto texto texto texto texto texto texto texto texto texto texto texto texto texto texto texto texto texto texto texto texto texto texto texto texto texto texto texto.

\section{Definição do Problema}

Texto texto texto texto texto texto texto texto texto texto texto texto texto texto texto texto texto texto texto texto texto texto texto texto texto texto texto texto texto texto texto texto texto texto texto texto texto texto texto \cite{ref:hodel}.

\section{Objetivo Geral}

Texto texto texto texto texto texto texto texto texto texto texto texto texto texto texto texto texto texto texto texto texto texto texto texto texto texto texto texto texto texto texto texto texto texto texto texto texto texto texto.

\section{Objetivos Específicos}

Texto texto texto texto:
\begin{alineas}
    \item texto texto texto texto texto texto texto texto texto texto texto texto texto texto texto texto texto texto texto texto texto texto texto texto texto texto texto texto texto texto texto texto texto texto texto texto texto texto;
    \item texto texto texto texto texto texto texto texto texto texto texto texto texto texto texto texto texto texto texto texto texto texto texto texto texto texto texto texto texto texto texto texto texto texto texto texto texto texto; 
    \item texto texto texto texto texto texto texto texto texto texto texto texto texto texto texto texto texto texto texto texto texto texto texto texto texto texto texto texto texto texto texto texto texto texto texto texto texto texto.
\end{alineas}
    
\section{Estrutura do Trabalho}

Texto texto texto texto texto texto texto texto texto texto texto texto texto texto texto texto texto texto texto texto texto texto. 
